\begin{frame}
    \frametitle{Prerequisites}
    \begin{itemize}
        \item Install a \LaTeX{} distribution
        \begin{itemize}
            \item \TeX{} Live % https://tug.org/texlive/
            \item \hologo{MiKTeX}
            \item Mac\TeX{}
        \end{itemize}
        \item Or use \href{https://www.overleaf.com/}{Overleaf} (account required!)
        \item Basic Syntax knowledge
    \end{itemize}
\end{frame}

\begin{frame}
    \frametitle{Syntax Basics}
    \framesubtitle{Commands}
    % [Speakers Notes]
    % Let's start with some basic syntax knowledge. 
    % One essential component of \LaTeX{} are \emph{commands} -- or sometimes also referred to as \emph{macros}.
    % Here's an example of a command, namely the `\section` command. 
    % It starts with a backslash -- a reserved character, but more no that later -- which tells \LaTeX{} that the following sequence of consecutive letters shall be interpreted as a command.
    % Those following letters are the commands name, which is then followed by a mandatory argument `{...}`. 
    % This argument is often some content, such as text.
    % In some cases there is also an optional argument as we see here `[...]`, 
    % though as the name implies, we can almost always omit this. 
    % These are the very basics of commands and we'll expand on this later on in the advanced section.
    \note{%
    \begin{itemize}
        \item Start with basic syntax knowledge
        \item Essential component of \LaTeX{} are \emph{commands} -- sometimes also referred to as \emph{macros}
        \item Here's an example
        \item Consists of
        \begin{itemize}
            \item Backslash -- mention briefly that it is a reserved character
            \item Command name (letters only)
            \item Argument
            \item Optional argument
        \end{itemize}
    \end{itemize}}
    \begin{center}
        % TODO: swap command with optional argument as last one to appear.
        \LARGE
        \tn{bslash}{\texttt{\bslash}}\tn{cmdname}{\textcolor{keywordcolour}{\texttt{\textbf{section}}}}\tn{cmdmand}{\texttt{\bropen A Section Heading\brclose}}
        
        \vspace{6em}
        
        \action<6->{%
            \LARGE
            \tn{bslash1}{\texttt{\bslash}}\tn{cmdname1}{\textcolor{keywordcolour}{\texttt{\textbf{section}}}}\tn{cmdopt}{\texttt{[A Heading]}}\tn{cmdmand1}{\texttt{\bropen A Section Heading\brclose}}
        }
    \end{center}
    \begin{tikzpicture}[%
        remember picture,
        overlay,
        thick,
        every node/.style={%
            text height=0.9em,
            text width=7em,
            align=center,
        }]
        \coordinate[below=0.5mm of bslash.south] (brackbase);
        \coordinate[below=2mm of bslash.south] (bracklow);
        \coordinate[above=0.5mm of cmdopt.north] (brackbasebot);
        \coordinate[above=2mm of cmdopt.north] (bracklowbot);
        \pause
        \node[below=7.5mm of bslash]  (bslashtxt) {Command\\ delimiter};
        \draw (brackbase-|{$(bslash.south west)+(0.5mm,0)$}) |- (bracklow-|bslash.south) -| (brackbase-|{$(bslash.south east)+(-0.5mm,0)$});
        \draw (bracklow-|bslash.south) -- (bslashtxt.north);
        \pause
        \node[right=0mm of bslashtxt] (cmdnametxt) {Command\\ name};
        \draw (brackbase-|{$(cmdname.south west)+(0.5mm,0)$}) |- (bracklow-|cmdname.south) -| (brackbase-|{$(cmdname.south east)+(-0.5mm,0)$});
        \draw (bracklow-|cmdname.south) |- ($(cmdname.south)!0.8!(cmdnametxt.north)$) -| (cmdnametxt.north);
        \pause
        \node[right=0mm of cmdnametxt]  (cmdopttxtphant) {\phantom{Optional}\\ \phantom{argument}};
        \node[right=0mm of cmdopttxtphant] (cmdmandtxt) {Mandatory\\ argument};
        \draw (brackbase-|{$(cmdmand.south west)+(0.5mm,0)$}) |- (bracklow-|cmdmand.south) -| (brackbase-|{$(cmdmand.south east)+(-0.5mm,0)$});
        \draw (bracklow-|cmdmand.south) |- ($(cmdmand.south)!0.8!(cmdmandtxt.north)$) -| (cmdmandtxt.north);
        \pause
        \node[below=-2mm of cmdopttxtphant] (cmdopttxt) {Optional\\ argument};
        \pause
        \draw (brackbasebot-|{$(cmdopt.north west)+(0.5mm,0)$}) |- (bracklowbot-|cmdopt.north) -| (brackbasebot-|{$(cmdopt.north east)+(-0.5mm,0)$});
        \draw (bracklowbot-|cmdopt.north) |- ($(cmdopt.north)!0.5!(cmdopttxt.south)$) -| (cmdopttxt.south);
    \end{tikzpicture}
\end{frame}

% NOTE: not sure if we still need the following...
% \begin{frame}
%     \frametitle{Syntax Basics}
%     \framesubtitle{Commands}
%     \begin{columns}
%         \begin{column}{.55\textwidth}
%             \begin{defblock}<1->
%                 \action<+->{\lstinline|\\commandName\{<argument>\}|\\}%
%                 \action<+->{\lstinline|\\commandName[<optional>]\{<argument>\}|}
%             \end{defblock}
%         \end{column}
%         \begin{column}{.44\textwidth}
%             \begin{exblock}<2->
%                 yadda yadda
%             \end{exblock}
%         \end{column}
%     \end{columns}
% \end{frame}

% TODO: rework this, it's currently too cluttered...
\begin{frame}
    \frametitle{Syntax Basics}
    \framesubtitle{Environments}
    % [Speakers Notes]
    % Next we are going to learn about environments. 
    %
    % Let's imagine we have some content -- i.e. some text -- that we want to center or 
    % generally just want to format differently. 
    % We could use a command, but once we have lots of content and different formatting, 
    % this gets quite convoluted and impractical. So we use environments!
    % An environment always surrounds its content with begin and end.
    % 
    % Here we have some content -- that ideally isn't just one line -- and we want to apply a
    % certain set of rules to all of it. For this, we can use so-called environments, which do exactly that.
    % For these environments there are two new commands that we need to know: `\begin{name}` and `\end{name}`.
    % These commands are a bit special, as they always have to come in pairs.
    % I.e. there is no `\begin{name}` without `\end{name}` and vice-versa.
    % The `name` can be any given predefined environment name -- in this case the name is `document`.
    \note{%
    Next up: environments
    \begin{itemize}
        \item Imagine we have content
        \item We want to apply certain rules to content
        \item Content ideally isn't just one line
        \item Could use commands, but gets convoluted and impractical with size
        \item Solution: environment surrounding content with begin and end
    \end{itemize}}
    \begin{center}
        \LARGE
        \parbox{20ex}{%
        \raggedright
        \onslide<2->{%
            \tn{envbegin}{\texttt{\bslash\lstkw{begin}}\texttt{\bropen}\tn{envname1}{\alt<3->{\texttt{center}}{\phantom{\texttt{center}}}}\texttt{\brclose}}\\%
        }
        \tn{envcontent}{\quad\textit{Content}}\\%
        \onslide<2->{%
            \tn{envend}{\texttt{\bslash\lstkw{end}}\texttt{\bropen}\tn{envname2}{\alt<3->{\texttt{center}}{\phantom{\texttt{center}}}}\texttt{\brclose}}%
        }
        }%
    \end{center}
    \begin{tikzpicture}[overlay, remember picture, thick]
        
    \end{tikzpicture}
\end{frame}

% \begin{frame}
%     \frametitle{Syntax Basics}
%     \framesubtitle{Environments}
%     \begin{center}
%         \LARGE
%         \parbox{20ex}{%
%         \raggedright
%         \tn{envbegin}{\texttt{\bslash\lstkw{begin}}\texttt{\bropen}\tn{envname1}{\texttt{document}}\texttt{\brclose}}\\%
%         \tn{envcontent}{\quad\texttt{...}}\\%
%         \tn{envend}{\texttt{\bslash\lstkw{end}}\texttt{\bropen}\tn{envname2}{\texttt{document}}\texttt{\brclose}}%
%         }%
%     \end{center}
%     \begin{tikzpicture}[overlay, remember picture, thick]
%         \coordinate[above=0.5mm of envbegin.north] (topbrackbase);
%         \coordinate[above=2mm of envbegin.north] (topbrackhigh);
%         \coordinate[below=0.5mm of envname1.south] (topbrackbaselow);
%         \coordinate[below=2mm of envname1.south] (topbracklow);
%         \coordinate[below=0.5mm of envend.south] (botbrackbase);
%         \coordinate[below=2mm of envend.south] (botbracklow);
%         \coordinate[above=0.5mm of envname2.north] (botbrackbasehigh);
%         \coordinate[above=2mm of envname2.north] (botbrackhigh);
%         \pause
%         \node[left=3cm of {$(envbegin)!0.5!(envend)$}, text width=12ex, align=center]  (delimtxt) {Environment delimiters};
%         \draw (topbrackbase-|{$(envbegin.north west)+(0.5mm,0)$}) |- (topbrackhigh-|envbegin.north) -| (topbrackbase-|{$(envbegin.north east)+(-0.5mm,0)$});
%         \draw (topbrackhigh-|envbegin.north) -- ++(0,1mm) -| (delimtxt);
        
%         \draw (botbrackbase-|{$(envend.south west)+(0.5mm,0)$}) |- (botbracklow-|envend.south) -| (botbrackbase-|{$(envend.south east)+(-0.5mm,0)$});
%         \draw (botbracklow-|envend.south) -- ++(0,-1mm) -| (delimtxt);

%         \node[right=3cm of {$(envname1)!0.5!(envname2)$}, text width=12ex, align=center]  (nametxt) {Environment name};
%         \draw (topbrackbaselow-|{$(envname1.south west)+(0.5mm,0)$}) |- (topbracklow-|envname1.south) -| (topbrackbaselow-|{$(envname1.south east)+(-0.5mm,0)$});
%         \draw (topbracklow-|envname1.south) -- (nametxt);
%         \draw (botbrackbasehigh-|{$(envname2.north west)+(0.5mm,0)$}) |- (botbrackhigh-|envname2.north) -| (botbrackbasehigh-|{$(envname2.north east)+(-0.5mm,0)$});
%         \draw (botbrackhigh-|envname2.north) -- (nametxt);
%     \end{tikzpicture}
% \end{frame}

% \begin{frame}
%     \frametitle{Syntax Basics}
%     \framesubtitle{Environments}
%     \begin{columns}
%         \begin{column}{.55\textwidth}
%                 \action<+->{%
%                 \lstinline|\\begin\{<environment-name>\}|\\
%                 \lstinline|...|\\
%                 \lstinline|\\end\{<environment-name>\}|}
%         \end{column}
%     \end{columns}
% \end{frame}

\begin{frame}
    \frametitle{Syntax Basics}
    \framesubtitle{Reserved Characters}
    % [Speakers Notes]
    % Now that we know commands and environments, 
    % there are also some reserved characters that we cannot type regularly. 
    % There are alternate ways to typeset them, so that we can use them in text if we need them.
    % As we've already seen before, the backslash is on here, but also the curly braces from before.
    % -------- Following might be better if we leave it out for now --------
    % In addition there's the underscore and and hat symbol, that are used in mathematical formulas as sub- and superscript.
    % There's also the dollar sign, which can be used to invoke the inline-math mode -- more on that later -- or the percentage symbol
    % which can be used for comments (similar to comments in programming languages, such as C, C++, Java, etc.).
    % The ampersand (&) is reserved for tabulators inside tables and the hash is used inside command definitions.
    \begin{center}
        \begin{tabular}{>{\ttfamily}c >{\ttfamily}l}
            \textnormal{\textbf{Symbol}} & \textnormal{\textbf{Substitute}}\\
            \toprule
            \bslash         & \bslash textbackslash\\
            \hline
            \bropen         & \bslash\bropen\\
            \hline
            \brclose        & \bslash\brclose\\
            \hline
            \char`\_        & \bslash textunderscore\\
            \hline
            \char`\^        & \bslash textasciicircum\\
            \hline
            \$              & \bslash\$\\
            \hline
            \%              & \bslash\%\\
            \hline
            \textasciitilde & \bslash textasciitilde\\
            \hline
            \&              & \bslash\&\\
            \hline
            \#              & \bslash\#\\
        \end{tabular}
        % \begin{tabular}{*{9}{>{\ttfamily}c}}
        %     \bslash & \bropen & \brclose & \char`\_ & \char`\^ & \$ & \% & \& &\# 
        % \end{tabular}
    \end{center}
\end{frame}

\begin{frame}
    \frametitle{Syntax Basics}
    \framesubtitle{Summary}
    \begin{columns}
        \begin{column}{.5\textwidth}
            % TODO: remove "Macro" from document as it may be confusing to have both terms. (maybe add as a footnote)
            \begin{block}{\blockstrut Commands}<1->
                \lstinline[keywords={commandName}]|\\commandName[<optional>]\{<argument>\}|\\%
                \lstinline[keywords={commandName}]|\\commandName\{<argument>\}|
            \end{block}
            \begin{block}{\blockstrut Environments}<2->
                \lstinline|\\begin\{<environment-name>\}|\\%
                \lstinline|  ...|\\%
                \lstinline|\\end\{<environment-name>\}|
            \end{block}
        \end{column}
        \begin{column}{.47\textwidth}
            \begin{block}{\blockstrut Reserved Characters}<3->
                \begin{tabular}{*{10}{>{\ttfamily}c}}
                    \bslash & \bropen & \brclose & \char`\_ & \char`\^ & \$ & \% & \textasciitilde & \& &\# 
                \end{tabular}
            \end{block}
        \end{column}
    \end{columns}
\end{frame}